% Transcription — fonds Alexandre Grothendieck, shelfmark 29, batch 6 (pages 101-120).
% Established from the facsimile archives/batches/29/batch-06.pdf (a mirror of
% https://grothendieck.umontpellier.fr/29.pdf), per the skill
% transcribe-grothendieck.
%
% Seventeen of the twenty pages are transcribed. Three are not:
%
%   - page 108, a sheet of a polemical typescript on Bourbaki ("Ad majorem
%     functorem gloriam"), scanned upside down because page 107 is written on
%     its back. It carries no mathematics;
%   - page 111, an empty orange folder cover inscribed, in his hand, "Notes
%     pour SGA 7". It is the wrapper of the typescript that follows, and names
%     it; as a separator sheet it gets no \page{};
%   - page 117, the blank verso of page 116, showing only that page's
%     bleed-through and one struck fragment — a formula number (2.1.14.6) and
%     four barred letters — that no reading can make sense of.
%
% Page 105 is kept although it is an upside-down verso of an unrelated text.
% It is a sheet of the EGA IV typescript, chapter IV page 878, running from
% Proposition (17.4.10) into section 17.5 and Théorème (17.5.1) on smooth
% morphisms — and it carries substantive corrections in his hand: "noethérien"
% struck and replaced by "de type fini", the two equivalences labelled in the
% margin, and the warning "Attention à l'ambiguïté dans « lisse » !". That is
% the criterion batch 5 used for its pages 82 and 84: a typed verso is kept
% when it carries something in his hand, and skipped when it does not.
%
% The batch falls into five runs.
%
% Page 101 is a carbon of a French typescript on local symbols on an abelian
% variety — the Albanese kernel Z_l(A), the symbol j_o and its invariance under
% translation, and the reciprocity i(X_a, b) = i(X_b^-, a) — with one insertion
% in his hand.
%
% Pages 102 to 104 and 106 are four unpaginated manuscript sheets, each on its
% own question. Page 102, numbered 3 by the author in a circle, redrafts as a
% corollary the Kummer-theoretic programme of page 100 of batch 5: the class of
% the extension 1 -> I -> G -> Gamma -> 1 in H^2, its image in H^2(D, ...), and
% the Chern class of the normal bundle. It is written very fast and struck
% through repeatedly; a great deal of it is left \ill{}, which is the honest
% reading of a 228 ppi scan of that hand and not a shortcut. Page 103, crossed
% out by two long diagonal strokes, is the quasi-unipotence lemma — F G F^{-1}
% = G^q forces the eigenvalues of G to be roots of unity — and it is the one
% page of the batch that reads cleanly from end to end. Page 104 is three lines
% on pi_2 and a perspective sketch of three divisors crossing. Page 106 draws
% two branches of a curve through a point, the local scheme Z = Spec O^_{X,a}
% around the crossing, and the diagram of fundamental groups the descent
% theorem would have to close.
%
% Pages 107, 109 and 110 are three sheets of commutative algebra: factoring a
% monic polynomial P = QQ' against decompositions of A[t]/(P), the criterion
% Q + Q' = A <=> Q ∩ Q' = QQ', and a sketch of the completion A -> A' -> A^.
%
% Pages 112 to 115 are the typescript "Théorème de Bertini et variantes", with
% his corrections: Bertini 1, the local statement 2, corollaries 3, 4, 5 on the
% surjectivity of pi_1(Y_t) -> pi_1(X), theorem 6 for the tame pi_1, and a
% remark 7 that stops on what he cannot prove for a general pencil.
%
% Pages 116, and 118 to 120, are the typescript "Comportement fonctoriel des
% groupes d'inertie", sections 1 to 5, likewise corrected in his hand.
%
% The typescripts are transcribed as he left them: the typist's xxxx-ed out
% words and his own deletions as \struck{}, his handwritten insertions as
% \add{}, his marginal notes as \marginal{}. Where a typed word is struck and a
% handwritten one put above it, both are given, struck first.
%
% The dating is the inventory's own, for the whole shelfmark.
%
% Pass: Opus 5 (claude-opus-5), 2026-08-25 — first pass, unchecked against the pages by a human.
\documentclass[11pt,a4paper]{article}
% Le préambule commun à toutes les transcriptions du fonds Grothendieck.
%
% Il tient en une idée : **l'incertitude se compose, elle ne se résout pas.**
% Une transcription qui lisse un mot douteux a détruit la seule chose pour
% laquelle on la faisait — un lecteur doit pouvoir savoir, à chaque endroit,
% ce qui était sur le feuillet et ce qui a été ajouté. Les six macros qui
% suivent sont donc l'appareil critique, et non de la décoration.
%
% Les mêmes commandes sont reconnues par `scripts/render.mjs`, qui produit la
% vue de lecture du site. Une macro ajoutée ici doit l'être aussi là-bas,
% faute de quoi le rendu échouera bruyamment — ce qui est voulu : un rendu
% silencieusement faux est pire qu'un rendu absent.

\ProvidesPackage{grothendieck}[2026/08/08 Transcriptions du fonds Grothendieck]

\usepackage{amsmath,amssymb,amsthm}
\usepackage{mathrsfs}
\usepackage{stmaryrd}
% Les diagrammes commutatifs sont partout dans ce fonds : c'est la langue
% même de la géométrie algébrique de Grothendieck, et une transcription qui
% les rendrait en prose ne transcrirait rien.
\usepackage{tikz-cd}
% Français par défaut — c'est la langue des feuillets — mais l'anglais est
% chargé aussi : la traduction et le résumé sont des documents anglais, et la
% typographie française (espaces avant les deux-points) y serait une faute.
% Ces documents basculent par \selectlanguage{english} en tête de corps.
\usepackage[english,french]{babel}
\usepackage{fontspec}
\usepackage{geometry}
\geometry{a4paper,margin=25mm,marginparwidth=28mm}
\usepackage{marginnote}
\usepackage{xcolor}
% ulem, not soul. `soul`'s \st and \ul are built on a character-by-character
% scan that XeLaTeX's font machinery breaks: the document fails with an
% undefined control sequence rather than degrading. ulem does the same two jobs
% and is Unicode-engine safe. `normalem` stops it hijacking \emph, which the
% transcriptions use for Grothendieck's own emphasis.
\usepackage[normalem]{ulem}
\usepackage{hyperref}
\hypersetup{colorlinks=true,linkcolor=black,urlcolor=black,pdfborder={0 0 0}}

% --- Métadonnées du lot -----------------------------------------------------
% Reprises telles quelles par le script de rendu, qui les affiche en tête de
% la vue de lecture. Elles fixent ce que le fichier prétend couvrir : sans
% elles, une transcription flotte sans point d'attache dans un fonds de seize
% mille feuillets.

\newcommand{\folder}[1]{\gdef\grothFolder{#1}}
\newcommand{\batch}[1]{\gdef\grothBatch{#1}}
\newcommand{\leaves}[2]{\gdef\grothFirst{#1}\gdef\grothLast{#2}}
\newcommand{\foldertitle}[1]{\gdef\grothFoldertitle{#1}}
\gdef\grothFolder{?}\gdef\grothBatch{?}\gdef\grothFirst{?}\gdef\grothLast{?}\gdef\grothFoldertitle{}

% --- Le résumé --------------------------------------------------------------
%
% Il ouvre la lecture modernisée, sous le titre « Résumé », et il est composé
% en petit. Ce n'est pas un ornement : c'est le seul passage du document qui
% ne soit pas des mathématiques, et un lecteur doit pouvoir le reconnaître
% avant de l'avoir lu — puis le sauter s'il connaît déjà le sujet, ou s'y
% arrêter s'il ne le connaît pas. Le filet qui le suit marque la reprise du
% corps.

\newenvironment{resume}
  {\small\setlength{\parskip}{0.45em}}
  {\par\medskip\hrule\medskip}

% --- Mots-clés ---------------------------------------------------------------
%
% En anglais, en clôture du résumé de la lecture modernisée : le vocabulaire
% moderne sous lequel chercher ce que les feuillets construisent. Le manifeste
% du site (`npm run manifest`) les extrait pour étiqueter la cote entière —
% cette ligne est la source unique des tags, il n'y a pas de fichier à côté.
\newcommand{\keywords}[1]{\par\smallskip\noindent
  {\footnotesize\textsc{Keywords} --- \textit{#1}}\par}

% --- L'appareil critique ----------------------------------------------------

% Le numéro de feuillet, en marge. C'est l'ancre qui permet de régler un
% désaccord sur une formule en regardant *un* feuillet plutôt qu'une chemise
% de six cents. Le site s'en sert aussi pour tourner les pages du fac-similé
% quand on fait défiler la transcription.
\newcommand{\leaf}[1]{%
  \marginnote{\footnotesize\color{gray!70}\textsf{#1}}%
  \label{leaf:#1}\ignorespaces}

% Illisible : on ne devine pas. Un mot inventé qui se lit comme les autres est
% le pire résultat possible.
\newcommand{\ill}{\textcolor{red!60!black}{[\,\ldots\,]}}

% Lecture douteuse : le mot est donné, et signalé comme incertain.
\newcommand{\uncertain}[1]{\uline{#1}}

% Ajout de l'éditeur — une lettre manquante, un mot sous-entendu.
\newcommand{\add}[1]{\textcolor{blue!50!black}{[#1]}}

% Biffé par Grothendieck lui-même. Ce qu'il a rayé fait partie du document :
% une rature dit souvent où la pensée a changé de direction.
\newcommand{\struck}[1]{\sout{#1}}

% Note du transcripteur, sur l'état matériel du feuillet.
\newcommand{\note}[1]{\par\smallskip{\small\color{gray!60!black}\textit{[#1]}}\par\smallskip}

% Note marginale de Grothendieck, distincte du corps du texte.
\newcommand{\marginal}[1]{\par\smallskip\noindent\hspace{1em}%
  {\small\color{gray!60!black}#1}\par\smallskip}

% --- Titre ------------------------------------------------------------------

\AtBeginDocument{%
  \begin{center}
    {\large\bfseries Fonds Alexandre Grothendieck --- cote n\textsuperscript{o}~\grothFolder}\\[2pt]
    {\itshape\grothFoldertitle}\\[4pt]
    {\small feuillets \grothFirst--\grothLast\ (lot \grothBatch)}\\[2pt]
    {\footnotesize\color{gray!60!black}Transcription établie d'après le fac-similé numérisé
     par l'Université de Montpellier.\\
     Les lectures incertaines sont soulignées, les illisibles notées [\,\ldots\,],
     les ajouts entre crochets.}
  \end{center}
  \vspace{1em}}

\endinput


\folder{29}
\batch{6}
\pages{101}{120}
\dating{1959-1969}
\watermark{Édition de démonstration}
\foldertitle{Groupe fondamental [Autour de SGA 4 et SGA 7] : lettres (1959, 1967, 1969), tapuscrits et copies de tapuscrit annotés (s.d.), notes manuscrites (s.d.)}

\begin{document}

\section*{Symboles locaux : le cas du noyau d'Albanese (page 101)}

\page{101}

Examinons maintenant le cas où l'un des cycles $\underline{a}$ ou
$\underline{b}$ (par exemple $\underline{b}$) appartient au noyau d'Albanese
$Z_{\ell}(A)$ ; par linéarité, on se ramène au cas où $\underline{b}$ est de la
forme $\underline{c}_{u} - \underline{c}$, avec $\underline{c} \in
Z_{o}^{s}(A)_{\underline{k}}$, et $u \in {}^{s}\!A_{\underline{k}}$. D'après
l'invariance du symbole $j_{o}$ par la translation $\tau_{u}$, les deux membres
de (3) sont respectivement égaux à $j_{o}(Y_{\underline{a}}, \underline{c})$ et
à $j_{o}(Y^{-}_{\underline{c}}, \underline{a})$, \add{en posant $Y = X_{u} -
X$.} Puisqu'on a $X \in D_{\underline{a}}(A)$, on a $Y \sim 0$, et on est ramené
au cas déjà traité.

Supposons maintenant $X$, $\underline{a}$ et $\underline{b}$ quelconques. On
peut, en modifiant $\underline{b}$ par l'addition d'un élément de
$Z_{\ell}^{s}(A)_{\underline{k}}$, se ramener au cas où le cycle $\underline{a}^{-}
\otimes \underline{b}$ est étranger à $X$. Les symboles
$X_{\underline{a}}(\underline{b})$ et $X^{-}_{\underline{b}}(\underline{a})$
sont alors définis, et on a

\[ X_{\underline{a}}(\underline{b}) = X^{-}_{\underline{b}}(\underline{a}),
\qquad \text{d'où} \qquad v\big(X_{\underline{a}}(\underline{b})\big) =
v\big(X^{-}_{\underline{b}}(\underline{a})\big). \]

D'autre part, soit $a$ (resp.\ $b$) l'un quelconque des composants de
$\underline{a}$ (resp.\ $\underline{b}$). Alors les symboles $i(X_{a}, b)$ et
$i(X^{-}_{b}, a)$ sont définis ; l'application $\zeta : A \to A$ définie par
$\zeta(x) = a + b - x$ étant $\underline{p}$-isomorphique en $b^{o}$, de valeur
$a^{o}$, on a, d'après la propriété d'invariance du symbole $i$,

\[ i(X_{a}, b) = i(X^{-}_{b}, a). \]

Il en résulte, par linéarité,

\[ i(X_{\underline{a}}, \underline{b}) = i(X^{-}_{\underline{b}},
\underline{a}). \]

Compte tenu des formules (2) précédentes, on en déduit (3).

Notons enfin que $j_{o}(X, \underline{a})$ s'annule lorsque tous les composants
de $A$ appartiennent au groupe ${}^{s}\!A_{\underline{k}, o}$. En effet, soit
$f$ une fonction sur $V$, définie sur $\underline{k}$, telle que

\section*{Corollaire kummérien pour un revêtement modérément ramifié (page 102)}

\page{102}
\note{Feuillet numéroté ③ par l'auteur. Il est écrit très vite et raturé d'un
bout à l'autre : ce qui est porté ici est ce qui se lit, le reste est laissé
illisible. Il reprend, sous forme de corollaire, le programme kummérien
esquissé page 100 du carnet 5.}

\marginal{$K \;\; \tilde{K} \;\; K'$ au-dessus de $X \;\; \tilde{X} \;\; X'$}

\textbf{Corollaire.} $X$ schéma normal \struck{et} loc.\ intègre, $D_{i}$
diviseurs \add{de Cartier} \struck{irréductibles} \ill{}, $X'$ revêtement
\struck{galoisien} \ill{} \ill{} les groupes d'inertie \ill{} relatifs \ill{}
invariants \ill{} \ill{} $X'$ soit modérément \ill{} $D_{i}$. Donc \ill{}

\[ 1 \longrightarrow I \longrightarrow G \longrightarrow \Gamma \longrightarrow
1 \]

\struck{$I$ est le produit des \ill{}} De plus, $X'/I = \tilde{X}$ et \ill{} un
revêtement principal de $X$ de groupe $\Gamma$. Donc \ill{} donne une classe

\[ \alpha \in H^{2}(\Gamma, I) \]

\ill{} de l'inertie \ill{} définit une classe

\[ \alpha \in H^{2}(X, \tilde{X} \times^{\Gamma} I), \]

\struck{d'où \ill{}} On a \ill{} canonique \ill{} $H^{2}(D, \tilde{D}
\times^{\Gamma} I)$. On a \ill{} isomorphisme canonique

\[ \tilde{D} \times^{\Gamma} I \;\simeq\; \prod_{i} \mu_{n_{i}, D} \quad
\text{sur } D \]

\ill{} compatible avec les opérations de $\Gamma$ \ill{} sur $(\tilde{X},
\Gamma)$. \ill{}

\[ I_{\tilde{X}} \simeq \mu_{\underline{n}} \]

relatif aux $\tilde{D}_{i}$ ; comparer $\underline{n} = (n_{i})$. \ill{}
$\alpha$ dans $H^{1}(D, \mu_{\underline{n}})$ \ill{} intègre. Le choix d'un
\ill{}

\[ I_{i} \simeq \mu, \qquad c\big(\mathcal{O}_{X}(D)\big) \in H^{2}(X,
\mu_{\underline{n}}), \]

i.e.

\[ c\big(T_{D/X}\big) \in H^{2}(D, \mu_{\underline{n}}), \quad \text{fibré
normal.} \]

\section*{Quasi-unipotence de la monodromie locale (page 103)}

\page{103}
\note{La page est barrée de deux longs traits diagonaux.}

\ill{} (par le choix d'isomorphismes $Z_{\ell}(1) \simeq Z_{\ell}$ pour tout
$\ell$) \ill{} on définit une représentation \ill{} le produit semi-direct de
$\mathbb{Z}$ avec $Z_{\ell}$, opérant sur le sous-groupe $Z_{\ell}$ par
multiplication par $q$. Donc ce groupe est engendré par deux éléments $f$
(correspondant à $f\circ b_{\rho}$) et $g$, satisfaisant la relation

\begin{equation*}
(6.4) \qquad f g f^{-1} = g^{q}
\end{equation*}

\textbf{Lemme 6.5.} Soient $F$, $G$ deux \struck{\ill{}} endomorphismes d'un
espace vectoriel $E$ de dimension finie sur un corps, et $q$ un entier $> 1$
tels qu'on ait \struck{(6.4)}

\[ F G F^{-1} = G^{q}. \]

Alors $G$ est quasi-unipotent, i.e.\ ses valeurs propres sont toutes des racines
de $1$, ou encore : il existe un entier $N \geqslant 1$ tel que $G^{N}$ soit
unipotent.

\emph{Dém.} \struck{Les} Si $G$ et \ill{} dans une clôture alg.\ $\bar{R}$ de
$R$) \struck{étant} les \struck{mêmes}, il \ill{} $=$ celles de $F G F^{-1}$,
donc \ill{} celles de $G^{q}$, il s'ensuit que c'est une partie $\Phi$ de
$\bar{R}^{*}$ stable par l'application

\[ \lambda \longmapsto \lambda^{q}. \]

Comme $\Phi$ est fini, il existe donc pour tout $\lambda \in \Phi$ deux entiers
$N' \neq N''$ \struck{tels} $\geqslant 0$ tels que \ill{}

\section*{$\pi_{2}$ et croisements normaux (page 104)}

\page{104}

$\pi_{2}(X, \xi)$

\begin{enumerate}
\item[a)] Si $X$ simplement connexe, alors $M \rightsquigarrow H^{2}(X, M)$,
avec $M$ \struck{\ill{}} groupe abélien fini, est exact ; donc
pro-représentable par
\[ H^{2}(X, M) = \operatorname{Hom}\nolimits_{\mathrm{cont}}\big(\pi_{2}(X),
M\big). \]

\item[b)] Dans le cas général, $\pi_{2}(X, \xi) = \pi_{2}\big(\tilde{X}(\xi)\big)$,
$\tilde{X}(\xi)$ revêtement universel pointé défini par $\xi$.
\end{enumerate}

\note{Le bas de la page porte un croquis en perspective, au crayon, de trois
diviseurs $D_{0}$, $D_{1}$, $D_{2}$ se coupant à croisements normaux : trois
plans deux à deux sécants, leurs trois droites d'intersection concourantes en
un point. Il n'est pas repris en diagramme, n'étant pas commutatif.}

\section*{Caractérisations des morphismes lisses : feuillet d'EGA IV annoté (page 105)}

\page{105}
\note{Page IV-878 du tapuscrit d'EGA IV, portant les corrections de l'auteur.}

un point de $Y$.

Il résulte en effet de $(17.4.1,(b''))$ qu'une $Y$-section $s$ de $X$ est une
immersion ouverte, et comme $X$ est un $Y$-schéma, $s$ est aussi une
\struck{\ill{}} immersion fermée $(I, 5.4.7)$ ; il en résulte que $s$ est un
isomorphisme de $Y$ sur un sous-préschéma de $X$ induit sur une partie ouverte
et fermée de $X$, et comme $s(Y)$ est connexe, c'est nécessairement une
composante connexe de $X$. Le reste de la proposition est immédiat.

\textbf{Proposition (17.4.10).} --- \emph{Soient $g : Y \to S$, $h : X \to S$
deux morphismes localement de présentation finie. Pour qu'un $S$-morphisme $f :
X \to Y$ soit non ramifié il faut et il suffit que pour tout $s \in S$, le
morphisme $f_{s} : h^{-1}(s) \to g^{-1}(s)$ déduit de $f$ par le changement de
base $\operatorname{Spec}(\kappa(s)) \to S$, soit non ramifié.}

Cela résulte aussitôt du fait que pour un morphisme localement de présentation
finie, le fait d'être non ramifié est une propriété des fibres de ce morphisme
\add{$(17.4.1, d))$}, et que pour tout $y \in Y$, on a $f^{-1}(y) =
f_{s}^{-1}(y)$ si $s = g(y)$.

\subsection*{17.5. Caractérisations des morphismes lisses.}

\textbf{Théorème (17.5.1).} --- \emph{Soient $f : X \to Y$ un morphisme
localement de présentation finie, $x$ un point de $X$, $y = f(x)$. Les
conditions suivantes sont équivalentes :}

\begin{enumerate}
\item[a)] \emph{$f$ est lisse au point $x$.}
\item[b)] \emph{$f$ est plat au point $x$ et le $\kappa(y)$-préschéma
$f^{-1}(y)$ est lisse sur $\kappa(y)$ au point $x$.}
\item[b')] \emph{$f$ est régulier au point $x$ $(6.8.1)$.}
\item[c)] \emph{L'anneau $\mathcal{O}_{X,x}$ est une
$\mathcal{O}_{Y,y}$-algèbre formellement lisse pour les topologies discrètes.}
\end{enumerate}

On peut se borner au cas où $Y = \operatorname{Spec}(A)$, $X =
\operatorname{Spec}(C)$, où $C = B/J$, $B = A[T_{1}, \dots, T_{n}]$ étant une
algèbre de polynômes et $J$ un idéal de $B$ de type fini. L'équivalence de a)
et c) résulte alors de $(\underline{0}, 22.6.4)$. \marginal{a) $\Leftrightarrow$
c)} D'autre part, comme $f^{-1}(y)$ est un $\kappa(y)$-préschéma localement
\struck{noethérien} \add{de type fini}, l'équivalence de b) et b') résulte de
$(6.8.6)$. \marginal{a) $\Leftrightarrow$ b)} Reste à voir que a) et b) sont
équivalentes.

\marginal{Attention à l'ambiguïté dans « lisse » !}

\section*{$\pi_{1}$ local le long de deux branches (page 106)}

\page{106}
\note{Le haut de la page est un dessin : deux branches, $X_{T'}$ horizontale et
$X_{T''}$ verticale, se coupant en un point, chacune hachurée en bandes
transverses ; autour du point de croisement, un petit disque pointillé porte la
lettre $Z$. C'est la figure que les égalités qui suivent décrivent.}

\[ Z = \operatorname{Spec} \hat{\mathcal{O}}_{X,a} \]

\[ X_{T'} \cap X_{T''} = p \]

$X_{T'} \cap Z = Z' = \hat{A}_{p'}$, \quad correspond au corps résiduel $k(p')
= $ corps local en $a$ de $T'$.

$X_{T''} \cap Z = Z'' = \hat{A}_{p''}$, \quad correspond au corps résiduel
$k(p'') = $ corps local en $a$ de $T''$.

\begin{tikzcd}[column sep=small, row sep=small, nodes={font=\scriptsize}]
 & \hat{Z} = \pi_{1}(Z') \arrow[dl] \arrow[dr] & \pi_{1}(Z'') = \hat{Z}
 \arrow[dl] \arrow[d] \\
\pi_{1}(X_{T'}) \arrow[dr] & \pi_{1}(X_{T''}) \arrow[d] & \pi_{1}(Z) =
\hat{Z}^{2} \arrow[dl] \\
 & \pi_{1}(X_{T}) &
\end{tikzcd}

On doit prouver un th.\ de descente \uncertain{de nature essentiellement
locale} $+$ déterminer certains hom.\ canoniques.

\section*{Décomposition d'un polynôme unitaire et idéaux étrangers (pages 107 à 110)}

\page{107}

\[ P = QQ', \qquad QR + Q'R' = 1, \qquad (P) = (Q)(Q') = (Q) \cap (Q') \]

\[ A[t]/(P) \longrightarrow A[t]/(Q) \times A[t]/(Q') \]

$s \in (Q)$, $s \in (Q') \Rightarrow s = Q'sR' + QsR \in (QQ')$.

Injectif car $(P) = (Q) \cap (Q')$ ; surjectif car \ill{} Nakayama \ill{}

\textbf{Lemme.} Soient $Q$, $Q' \subset A$ \struck{\ill{}} deux idéaux. Alors

\[ A/(QQ') \longrightarrow A/Q \times A/Q' \]

\struck{est un isomorphisme} \ill{}

\begin{enumerate}
\item[(i)] $Q \cap Q' = Q\,Q'$ ;
\item[(ii)] $Q + Q' = A$.
\end{enumerate}

\[ Q \cap Q' = Q\,Q'\,(Q + Q') \subset Q'Q + QQ' = QQ' \]

\ill{} $\forall\, x, y \in A$, \ill{} $x \pmod{Q}$, $y \pmod{Q'}$ \ill{}

\[ A/Q\,Q' \longrightarrow A/Q \times A/Q' \quad \text{surjectif} \]

\[ Q + Q' = A \iff QQ' = Q \cap Q' \]

\page{109}

\begin{enumerate}
\item[a)] Donnons-nous \struck{un} un idéal $R \subset B$. Alors décompositions
de $B/R$ en produit de deux anneaux $\simeq$ données de deux idéaux $J$, $J'$
de $A$ tels que
\[ J J' = R, \qquad J + J' = A \]
[car $J + J' = A \Rightarrow J J' = J \cap J'$]. \marginal{\ill{} $A$
commutatif}

\item[b)] Supposons $B = A[T]$, $(R = (P))$, $P$ polynôme unitaire. Alors les
décompositions de $R$ en produit de deux idéaux $J$, $J'$ tels que $J + J' = B$
sont \uncertain{uniques} de la forme
\[ (P) = (Q)\,(Q'), \]
i.e.\ $Q$, $Q'$ sont des polynômes unitaires tels que $P = QQ'$.
\end{enumerate}

Si l'on ne suppose pas une telle déc., alors $P$ \uncertain{est unique} à une
unité près, et \ill{} \ill{} unique \ill{} [\uncertain{petite vérification}
\ill{}]. \struck{\ill{}} De l'unicité résulte qu'on peut supposer
$\operatorname{Spec} A$ local. Alors pour l'existence, on note qu'il est vrai
sur le \ill{}

\page{110}

\struck{\ill{} par contre \ill{} pourrait \ill{} m'inquiéter sérieusement}

\[ A \text{ anneau local}, \qquad \hat{A} \text{ son complété} \]
\[ X = \operatorname{Spec}(A), \qquad \hat{X} \]
\[ D \text{ partie fermée}, \qquad \hat{D} \]

\[ \hat{Y} \hookrightarrow \hat{Z}, \qquad \hat{Y} \times_{\hat{S}} \hat{Z} \]

$Y \times Z$ \ill{} composante connexe \ill{} $Y \times Z$.

\begin{tikzcd}
S & \hat{S} \arrow[l] \\
X \arrow[u] & \hat{X} \arrow[l] \arrow[u]
\end{tikzcd}

\[ A \longrightarrow A' \longrightarrow \hat{A} = \hat{A}' \]

\note{Trois petits cubes en perspective, chacun marqué d'un point intérieur,
figurent le schéma local et son point fermé ; ils ne sont pas repris.}

\section*{Théorème de Bertini et variantes (pages 112 à 115)}

\page{112}

\subsection*{Théorème de Bertini et variantes.}

Dans la suite $k$ désigne un corps, $P^{r} = P^{r}_{k}$.

\textbf{Théorème de Bertini 1 :} Soit $f : X \to P^{r}_{k}$ un morphisme, où
$X$ est un schéma algébrique géométriquement irréductible, et soit $n = \dim
\overline{f(X)}$. Soit $L_{t}$ le \struck{\ill{}} sous-espace linéaire de
codimension $s$ \emph{générique} de $P^{r}$, où $s$ est un entier fixé
$\leqslant n-1$. Alors \add{$Y_{t} =$} $f_{t}^{-1}(L_{t}) \subset X_{t} =
X \otimes_{k} k(t)$ est géométriquement irréductible.

\marginal{\uncertain{vrai} dans EGA \uncertain{IV}}

\textbf{Énoncé local 2 :} Soit $f : X \to P^{r}$ un morphisme, avec $X$ un
schéma algébrique. \struck{Complémentairement supposons} Soit $Z \to X$
\struck{\ill{}} l'ensemble des points de $X$ en lesquels $f$ est \add{soit}
ramifié \struck{ou en lesquels} \add{soit} $X$ n'est pas géométriquement
unibranche (resp.\ \add{pas} normal, resp.\ \add{pas} lisse \dots),
\struck{alors} et supposons que $\overline{f(Z)}$ soit de dimension $< s$.
Alors \add{$Y_{t} =$} $f_{t}^{-1}(L_{t})$ est géométriquement unibranche
(resp.\ normal, resp.\ lisse \dots), où $L_{t}$ est comme ci-dessus.

De ces deux énoncés on déduit formellement ;

\textbf{\struck{Théorème} Corollaire 3 :} Soit $f : X \to P^{r}$ satisfaisant
les conditions de 1 et 2 \add{Supposons $k$ alg.\ clos.} non respé.\ Alors
\struck{l'inclusion} le morphisme canonique

\[ \pi_{1}(Y_{\bar{t}}) \longrightarrow \pi_{1}(X) \]

est surjectif.

\textbf{Remarque 4.} Utilisant les théorèmes de changement de base pour le
$\pi_{1}$, on conclut \struck{que} (posant $X_{\bar{t}} = X \otimes_{k}
k(\bar{t})$) que l'homomorphisme

\begin{equation*}
(*) \qquad \pi_{1}(Y_{\bar{t}}) \longrightarrow \pi_{1}(X_{\bar{t}})
\end{equation*}

est surjectif si $X$ est propre, et que l'homomorphisme

\begin{equation*}
(**) \qquad \pi_{1}^{!}(X_{\bar{t}}) \longrightarrow \pi_{1}^{!}(X_{\bar{t}})
\end{equation*}

est surjectif, où le $!$ désigne le plus grand quotient profini premier à
l'exposant car.\ de $k$. On n'affirme pas, cependant, que $(*)$ soit toujours
surjectif (sans hypothèse de propreté sur $X$).

On déduit, de 1, par une application facile du théorème de connexité :

\textbf{Corollaire 4 (de 1) :} Sous les conditions de 1) pour $f : X \to P^{r}$,
si $X$

\page{113}

est de plus propre, alors pour \emph{toute} sous-variété linéaire \add{$L_{t}$}
de codimension $s$, $Y_{t} = f_{t}^{-1}(L_{t})$ est géométriquement connexe.

Conjuguant avec 2, on trouve, par l'argument qui a permis de déduire 3 de 1 :

\textbf{Corollaire 5 :} \struck{Avec} \add{Sous} les hypothèses de 3,
\add{avec $X$ propre,} l'homomorphisme canonique

\[ \pi_{1}(Y_{\bar{t}}) \longrightarrow \pi_{1}(X) \; \big(\simeq
\pi_{1}(X_{\bar{t}}) \text{ par le théorème de changement de base}\big) \]

est surjectif.

\marginal{cf.\ SGA 1 X 2.11}

Nous désirons également un énoncé analogue à 5, donnant une surjectivité pour
des $\pi_{1}$ pour un $t$ non générique, mais sans hypothèse de propreté, et
pour des quotients plus gros du $\pi_{1}$ que ceux envisagés dans $(**)$. On
peut donner à ce sujet le résultat suivant :

\textbf{Théorème 6.} Supposons $k$ algébriquement clos, $X$ propre, $f : X \to
P^{r}$, \add{et} le sous-espace linéaire $L_{t}$ de $P^{r}$ de codimension $s$
(défini sur $k$, pour fixer les idées) \struck{\ill{}} et enfin le sous-schéma
fermé $D$ de $X$ tels qu'on ait ceci ; \struck{\ill{}}

\begin{enumerate}
\item[a)] \struck{\ill{}} $Y_{t} = f^{-1}(L_{t})$ est non vide.

\item[b)] Pour tout $x \in Y_{t}$, $X$ est lisse en $x$, $Y_{t}$ est lisse en
$x$, de codimension $s$ dans $X$ en $x$, $f$ est net en $x$.

\item[c)] Pour tout $x \in D_{t} = D \cap Y_{t}$, ou bien
$\operatorname{codim}_{x}(D_{t}, Y_{t}) \geqslant 2$, ou bien $D$ est un
diviseur sur $X$ en $x$ et induit sur $Y_{t}$ en $x$ un diviseur à croisements
normaux.
\end{enumerate}

Supposons $U = X - D$ irréductible, et considérons l'inclusion $U_{t} = U \cap
Y_{t} \to U$. \struck{\ill{}} Soit $\pi_{1}^{t}(U)$ le \struck{\ill{}} quotient
du groupe fondamental de $U$ \struck{\ill{}} qui classifie les revêtements de
$U$ qui sont modérément ramifiés en les composantes irréductibles de
codimension $1$ de $D$ dans $X$ qui rencontrent $Y_{t}$, \struck{\ill{}} et
$\pi_{1}^{t}(U_{t})$ le quotient \struck{analogue} de $\pi_{1}(U_{t})$ qui
classifie les revêtements qui sont modérément ramifiés en les composantes
irréductibles de codimension $1$ de $D_{t}$. On a un homomor-

\page{114}

phisme évident

\[ \pi_{1}^{t}(U_{t}) \longrightarrow \pi_{1}^{t}(U) \]

(un point-base géométrique de $U_{t}$ étant choisi). Cet homomorphisme est
\emph{surjectif}.

\emph{Démonstration.} Lorsqu'on remplace $L_{t}$ par $L_{t'}$ générique, alors
la conclusion est contenue dans le corollaire 3, en d'autres termes, si $V$ est
un revêtement étale \add{connexe} de $U$ \struck{\ill{}} (satisfaisant la
condition de modération envisagée --- en fait, ce n'est pas nécessaire, en
vertu de 3) alors sa restriction $V_{t'}$ au dessus de $U_{t'}$ est connexe. On
veut prouver que sa restriction $V_{t}$ au dessus de $U_{t}$ est
\add{également} connexe. Or pour un paramètre variable sur la grassmanienne des
$L_{t'}$ \add{dans un voisinage ouvert assez petit $T$ de $\bar{t}$},
\struck{et voisin de $t$, la famille} la famille des $Y_{t'}$ est un morphisme
propre et lisse $Y \to T$, et celle des $D_{t'}$ est un sous-schéma fermé
\struck{\ill{}} $\Delta$ de $Y$, qui est la réunion d'un sous-schéma
$\Delta_{1}$ qui est de codimension $\geqslant 2$ sur chaque fibre (et qui ne
compte pas, à cause du théorème de pureté relative) et d'un diviseur
$\Delta_{2}$ à croisements normaux relatifs. Le revêtement donné $V$ de $U$
définit un revêtement \struck{\ill{}} $V'$ de $U' = Y - \Delta$, et il faut
montrer que la restriction de celui-ci à la fibre géométrique de $t$ est
connexe, sachant qu'il en est ainsi de sa restriction à la fibre géométrique en
$t'$, la générisation maximale de $t$. Or par la pureté, on sait que $V'$ se
prolonge en un revêtement étale $V''$ de $U'' = Y - \Delta_{2}$, et par
hypothèse ce dernier est modérément ramifié sur les fibres, relativement au
diviseur $\Delta_{2}$. La conclusion est alors une conséquence du théorème
\struck{\ill{}} de spécialisation pour le groupe fondamental modéré, qui
figure(ra) dans SGA 1 XIII (exposé de Mme Raynaud).

\textbf{Remarque 7.} Avec les notations de SGA 7 XVIII, 6 s'applique pour tout
pinceau de Lefschetz $D$ pour fournir la surjectivité de \struck{quotients per-}

\page{115}

l'homomorphisme

\[ \pi_{1}^{t}(D - \check{D}) \longrightarrow \pi_{1}^{t}(P - \check{X}), \]

et donne le théorème de conjugaison des cycles évanescents (donc le théorème
d'irréductibilité de la cohomologie évanescente) pourvu qu'on ne soit pas dans
le cas car.\ $k = 2$ et $\dim X$ impaire (NB \struck{ce n'est pas} dans la
situation du th.\ de Griffiths, il n'y a donc pas d'ennui). Dans ce dernier
cas, par contre, je ne sais \struck{\ill{}} pas prouver \add{la conjugaison
sous $\pi_{1}(D - \check{D})$}, sauf pour un pinceau \emph{générique},
puisqu'alors la représentation envisagée est sauvage ; \add{j'ignore si elle
reste vraie pour un pinceau \emph{général}.} Noter cependant que la conjugaison
des sous-groupes de monodromie locale \emph{dans} $\operatorname{Aut}(H^{n-1}(Y))$
(qui est très élémentaire) reste vraie, bien entendu, et celle-ci suffit pour
prouver le théorème que j'ai proposé (au par.\ 6 de XVIII) caractérisant le cas
où la condition (A) n'est \emph{pas} vérifiée.

\section*{Comportement fonctoriel des groupes d'inertie (pages 116 à 120)}

\page{116}

\subsection*{Comportement fonctoriel des groupes d'inertie.}

\textbf{1.} Soient $X$ un schéma, $U$ un ouvert de $X$, $G$ un groupe profini,
$U'$ un revêtement principal de $U$ de groupe $G$, $x$ un point \add{géom.} de
$X$, $\bar{X}$ le localisé strict de $X$ en $x$, $\bar{U}$ l'image inverse de
$U$ dans $\bar{X}$, $\bar{U}' = U' \times_{U} \bar{U}$. On suppose donné de
plus un schéma $\bar{X}'$ sur $\bar{X}$, avec des opérations de $G$ sur
$\bar{X}'$, de telle façon que \struck{l'espace} \add{le $\bar{U}$-schéma} à
opérateurs $\bar{X}'|\bar{U}$ coïncide avec $\bar{U}'$.

Pratiquement, on sera intéressé par le cas où $x \in X - U$ et où $x$ est
adhérent à $U$, où $\bar{X}'$ est entier sur $\bar{X}$ et où
$\mathcal{O}_{\bar{X}} = \bar{p}(\mathcal{O}_{\bar{X}'})^{G}$, $\bar{p} :
\bar{X}' \to \bar{X}$ étant la projection canonique, --- de sorte que $G$ opère
transitivement sur l'ensemble fibre $\bar{X}'_{\bar{x}}$. De plus, on
\struck{\ill{}} sera \add{le plus souvent} sous des conditions où on peut
construire canoniquement $\bar{X}'$ en termes de $\bar{U}'$ ; par exemple, si
on suppose $X$ normal en $x$, donc $\bar{X}$ normal, et $\bar{U}$ non vide, on
définira $\bar{X}'$ comme le normalisé de $\bar{U}$ relativement à $\bar{U}'$
sur $\bar{U}$ (i.e.\ comme limite projective des normalisés dans les
$\bar{U}'(i)$ sur $\bar{U}$, $\bar{U}'$ étant écrit comme limite projective de
schémas principaux $U'(i)$ sur les groupes quotients finis $G(i)$ de $G$).

Si maintenant \struck{\ill{}} $\xi$ est un point de $\bar{X}'_{\bar{x}}$, on
peut considérer le groupe de stabilité $G_{\xi}$ de $\xi$, formé des $g \in G$
tels que $g.\xi = \xi$ ; ce groupe prendra le nom de \emph{groupe d'inertie
relatif à} $\xi$. Dans le cas envisagé plus haut où les points de $\bar{X}'$
sur $\bar{x}$ sont tous conjugués sous $G$, on voit que pour $\xi$ variable,
les $G_{\xi}$ parcourent une classe de sous-groupes conjugués de $G$, qu'on
appellera \add{(par abus de langage)} \emph{les sous-groupes d'inertie}
\struck{de} \emph{relatifs à} $x$. Noter qu'en toute rigueur cette classe de
sous-groupes \struck{\ill{}} n'est pas déterminée par la seule connaissance de
$x \in X$, mais dépend du choix du prolongement de $\bar{U}'$ en $\bar{X}'$.
Lorsque $X$ est normal en $x$ et que $x$ est spécialisation d'un point de $U$
(ce qui signifie que $x$ est adhérent à $U$ si $\operatorname{esp}(X)$ est
localement noethérien), il sera sous-entendu, sauf mention du contraire, que le
choix en question est fait par normalisation, comme expliqué plus haut.

\page{118}

\textbf{2.} Supposons donné une situation $(Y, V, H, V')$ analogue à $(X, U, G,
U')$, et un morphisme $f : X \to Y$ tel que $f(U) \to f(V)$, un morphisme $f' :
U' \to V'$ compatible avec $U \to V$ induit par $f$, un morphisme de groupes
$\phi : G \to H$ tel que \struck{\ill{}} $f'$ soit compatible avec les actions
de $G$ et $H$ via $\phi$.

Le cas le plus courant sera celui où on se donne un point géométrique $a$ de
$U$, \add{d'où} un point géométrique $b$ de $V$ \add{$= f(a)$}, où $U'$ et $V'$
sont respectivement des quotients des revêtements universels de $U$ et $V$
basés en $a$ et $b$, $G$ et $H$ étant respectivement des quotients de
$\pi_{1}(U, a)$ et $\pi_{1}(V, b)$, l'homomorphisme $f' : U' \to V'$ étant
induit par les homomorphismes canoniques sur les revêtements universels en $a$
resp.\ $b$, \struck{\ill{}} $\phi : G \to H$ étant induit par $f_{*} :
\pi_{1}(U, a) \to \pi_{1}(V, b)$ ; de cette façon $f'$ et $\phi$ sont
déterminés sans ambiguïté par $f$, et leur existence signifie simplement que
$f_{*} : \pi_{1}(U, a) \to \pi_{1}(V, b)$ veut bien passer au quotient en un
$\phi : G \to H$ (i.e.\ applique un certain sous-groupe dans un certain
sous-groupe).

Considérons d'autre part le localisé strict $\bar{Y}$ de $Y$ en le point
géométrique $f(\bar{x})$ de $Y$, d'où $\bar{V}$, $\bar{V}'$ comme dessus ;
supposons donné un prolongement $\bar{Y}'$ de $\bar{V}'$ en un schéma à groupe
d'opérateurs $H$ sur $\bar{Y}$, et supposons donné enfin un homomorphisme
$\bar{f}' : \bar{X}' \to \bar{Y}'$ qui prolonge l'homomorphisme $\bar{U}' \to
\bar{V}'$ induit par $f'$. Pour tout point $\xi$ de $\bar{X}'_{\bar{x}}$,
$\bar{f}'(\xi)$ est alors un point de $\bar{Y}'_{\bar{y}}$, et on peut
considérer le groupe d'inertie correspondant $H_{\bar{f}'(\xi)}$. On a alors la
relation évidente

\begin{equation*}
(2.1) \qquad \phi(G_{\xi}) \subset H_{\bar{f}'(\xi)}.
\end{equation*}

Dans le cas où le groupe $G$ resp.\ $H$ opère transitivement sur
$\bar{X}'_{\bar{x}}$ resp.\ $\bar{Y}'_{\bar{y}'}$, on peut donc dire que $\phi$
\struck{transforme la classe de conjugaison de sous-groupes} \add{applique tout
groupe} d'inertie de $G$ rel.\ à $x$ \add{dans un groupe d'}\struck{en la
classe de conj.\ des sous-groupes} d'inertie de $H$ rel.\ à $y = f(x)$.

\struck{D'autre part} Il reste alors à expliciter l'homomorphisme $\bar{f}' :
\bar{X}' \to \bar{Y}'$.

\page{119}

En pratique, il sera vrai que $\bar{U}'$ est schématiquement dense dans
$\bar{X}'$ \add{et $\bar{Y}'$ séparé sur $\bar{Y}$}, de sorte que $\bar{f}'$
sera déterminé de façon unique en termes des autres données, comme l'unique
prolongement de $\bar{U}' \to \bar{V}'$ induit par $f'$ ; ce sera le cas en
particulier dans le cas envisagé plus haut où $\bar{X}'$ se déduit de
\struck{\ill{}} $\bar{U}'$ par normalisation. Bien entendu, pour pouvoir
conclure \add{alors} que tout sous-groupe d'inertie de $G$ rel.\ à $x$ s'envoie
dans un sous-groupe d'inertie de $H$ rel.\ à $y$, il faudra avoir vérifié
l'\emph{existence} d'un prolongement $\bar{f}'$ de $\bar{U}' \to \bar{V}'$.

\add{(La commutation aux actions de $G$ sera alors \struck{automatique})}
\add{Cela \uncertain{marchera}, en interprétant $\bar{f}'$ comme $\bar{X}' \to
\bar{Y}' \times_{\bar{Y}} \bar{X}'$, lorsque $\bar{Y}'$ est entier
\struck{\ill{}} sur $\bar{Y}$, cf.\ EGA II 6.1.14.}

\textbf{3.} Revenons à la situation de 1), supposons que $\mathcal{O}_{X,x}$
soit un anneau local régulier, et que $\bar{U}$ soit le complémentaire dans le
schéma strictement local régulier $\bar{X}$ d'un diviseur \add{$\bar{D}$} à
croisements normaux. On est en particulier dans le cas favorable où la
définition de $\bar{X}'$ est immédiate par normalisation. Supposons que la
ramification de $\bar{U}'$ en les points maximaux de $\bar{D}$ soit
\emph{modérée}. On sait alors, par le calcul local bien connu, que pour tout
$\xi \in \bar{X}'_{\bar{x}}$, le groupe d'inertie correspondant $G_{\xi}$
s'identifie canoniquement \struck{\ill{}} à un groupe quotient du groupe

\begin{equation*}
(3.1) \qquad I = \prod_{\ell \neq p} \underline{Z}_{\ell}(1)\big(k(\bar{x})
\big)^{C} = \prod_{\ell \neq p} T_{\ell}\big(k(\bar{x})\big)^{C},
\end{equation*}

$p$ étant la caractéristique résiduelle en $x$.

\marginal{$C = $ l'ensemble des composantes irréd.\ de $\bar{D}$}

De plus, ce même calcul local montre que le facteur de ce groupe d'inertie
correspondant à la composante $\bar{D}_{c}$ ($c \in C$) du diviseur de
ramification $\bar{D}$ coïncide aussi avec le groupe d'inertie relatif à
n'importe quel point de $\bar{X}'$ au dessus du point maximal de $\bar{D}_{c}$
(ou de tout autre point de $\bar{D}_{c}$ qui n'appartient à aucune autre
composante irréductible de $\bar{D}$) qui générise $\xi$ ; comme il y a
toujours de tels points, cela montre que la classe de conjugaison du facteur en
question peut aussi s'interpréter comme la classe des \add{sous-}groupes
d'inertie en les points envisagés de $\bar{D}_{c}$ ; en particulier, cette
classe de conjugaison ne dépend que de la composante irréductible de $D = X - U$
\struck{\ill{}} de point générique l'image du point générique de $\bar{D}_{c}$.

\page{120}

\textbf{4.} Revenons à la situation de 2, et supposons que la structure de
$(\bar{X}, \bar{U})$ soit celle envisagée dans 3, et que de même $\bar{Y}$ soit
un schéma local régulier et $\bar{V} = \bar{Y} - \operatorname{supp} \bar{E}$,
$\bar{E}$ étant un diviseur à croisements normaux sur $\bar{Y}$. Soit $C'$
l'ensemble des composantes irréductibles de $\bar{E}$. On voit alors que pour
tout $c' \in C'$, l'image inverse de $\bar{E}_{c'}$ est un diviseur sur
$\bar{X}$ dont le support est contenu dans celui de $\bar{D}$, donc est de la
forme

\[ \bar{f}^{*}(\bar{E}_{c'}) = \sum_{c} n_{c,c'} \bar{D}_{c}, \]

où pour $c'$ fixé, les $n_{c,c'}$ sont des entiers $\geqslant 0$ non tous nuls.
Ces notations fixées, un calcul immédiat \struck{\ill{}} d'images inverses de
revêtements kummériens montre que pour deux points correspondants $\xi$ et $y$,
\struck{\ill{}} on a commutativité dans le diagramme

\begin{tikzcd}[column sep=small, row sep=small, nodes={font=\scriptsize}]
G_{\xi} \arrow[r, "\phi"] & H_{y} \\
\prod_{\ell} \underline{Z}_{\ell}(1)(\bar{x})^{C} \arrow[u] \arrow[r, "N"] &
\prod_{\ell} \underline{Z}_{\ell}(1)(\bar{y})^{C'} \arrow[u]
\end{tikzcd}

les flèches \struck{horizontales} \add{verticales} étant les épimorphismes
canoniques, et $N$ étant l'homomorphisme défini par la matrice $N =
(n_{c,c'})_{(c,c') \in C \times C'}$, compte tenu que l'homomorphisme $k(\bar{y})
\to k(\bar{x})$ induit des isomorphismes

\[ \underline{Z}_{\ell}(1)(\bar{y}) \xrightarrow{\;\sim\;}
\underline{Z}_{\ell}(1)(\bar{x}). \]

Pour que \add{l'homomorphisme} $N$ soit un isomorphisme, il est donc nécessaire
et suffisant que $\operatorname{card} C = \operatorname{card} C'$ et que
\struck{\ill{}} l'entier ordinaire $\det N$ soit égal à $\pm\, p^{r}$, où $p$
est l'exposant caractéristique de $k(x)$ et $k(y)$ ; dans ce cas, on conclut
donc que les \add{sous-}groupes d'inertie de $H$ en $y$ sont les conjugués des
images des \add{sous-}groupes d'inertie de $G$ en $x$.

\textbf{5.} Supposons en particulier que l'on ait $V = Y - \operatorname{supp}
E$, $E$ étant un diviseur \add{$\geqslant 0$} sur $Y$, \struck{\ill{}} $U =
f^{-1}(V)$, et que \struck{\ill{}} les conditions suivantes soient satisfaites :

\begin{enumerate}
\item[a)] \add{Pour $x \in f^{-1}(\operatorname{supp} E)$, $y = f(x)$,
$\mathcal{O}_{X,x}$ et $\mathcal{O}_{Y,y}$ sont noethériens.}
\item[b)] Le diviseur \add{$D =$} $f^{*}(E)$ est défini, c'est un sous-schéma
régulier de
\end{enumerate}

\end{document}
